\documentclass[10pt,a4paper]{amsart}
\usepackage[margin=19mm]{geometry}
\usepackage{amsmath,amssymb,mathtools}
\usepackage[scr=rsfs,cal=euler]{mathalfa}
\usepackage{xfrac}
\usepackage[unicode,hidelinks,bookmarks=false]{hyperref}
\newcommand{\ie}{i.e.\ }
\newcommand{\cf}{cf.\ }
\newcommand{\resp}{respectively\ }
\newcommand{\Subsection}{\subsection}
\providecommand{\nobreakdash}{\nobreak\hbox{-}}
\setlength{\emergencystretch}{2em}
\multlinegap0pt
\usepackage{amssymb}
\usepackage[all]{xy}
\usepackage{caption}
\usepackage{booktabs}
\usepackage{tikz}
\newtheorem{lemma}{Lemma}
\newtheorem{thm}{Theorem}
\renewcommand{\qedsymbol}{$\blacksquare$}
\title{On Bounds and Diophantine Properties of Elliptic Curves}
\author{Navvye Anand}
\date{Source: arXiv:2407.09558v1 (CC BY 4.0)}
\begin{document}
\maketitle

\noindent\emph{Source and licence.} On Bounds and Diophantine Properties of Elliptic Curves. Version arXiv:2407.09558v1 (CC BY 4.0), distributed by arXiv under the Creative Commons Attribution 4.0 International licence, http://creativecommons.org/licenses/by/4.0/, as recorded in the arXiv licence metadata for this exact version and shown on the abstract page https://arxiv.org/abs/2407.09558v1. Full licence notice: \texttt{licenses/arxiv-2407.09558-CC-BY-4.0.txt}.\par\smallskip

\noindent\textit{Editorial scope.} Counted unit: the article's thm thm (source0.tex lines 568--570). The article's complete original proof of it is delivered with it (source0.tex lines 571--593, one \texttt{proof} environment of 23 lines). No mathematics is written, repaired or re-derived: the statement and the proof are the article's own lines, in the article's own notation, and no retained line is substituted or re-flowed.  Retained uncounted as the source-local context the proof needs: lines 248--248; lines 544--546.  Nothing was omitted from any retained span, so the package carries no omission record.  No retained line contains a citation, so the wrapper carries no bibliography and no bibliography entry.  No retained line contains a figure environment, an image-inclusion command or drawing code, so no image is delivered and the package declares no asset dependency.  Editorial formatting only, in addition: an AMS \texttt{amsart} preamble that restates the article's own declarations and macros used by the retained lines, namely the environments \texttt{lemma}, \texttt{thm} on separate counters, together with the standard packages those lines need.  The retained environments print in this document as Lemma 1, Theorem 1 in order of appearance, whereas the article prints the same items with the numbering of its own full document.  The complete native source of the article as distributed in the arXiv e-print for this exact version is bundled unchanged as \texttt{sources/162246/source0.tex}.
\section{Binary Cubic Forms and Explicit Bounds for $N(E)$} \label{BinaryCubic}

\begin{lemma}[{Silverman's Theorem}]\label{Silverman's Theorem}
    Let $F$ be a binary cubic form with integer coefficients and non zero discriminant. Let $m_0 $ be a non-zero integer such that the curve $E : F(x,y) = m_0 z^3$ has a point over $\mathbb{Q}$. Using that point as origin, we give the $F$ the structure of an elliptic curve. Denote $r$ to be the Mordell-Weil rank of the elliptic curve over $\mathbb{Q}$, then there exists a constant $c$ which depends on $F$ such that $N_{F} > c \log{(m)}^{r/r+2}$ for infinitely many integers $m.$
\end{lemma}

\begin{thm}
Let $r$ be a positive integer which is the rank of a Mordell-Weil group of rational points of the elliptic curve $E: y^2 = x^3 + D$. Then, there exist infinitely many inequivalent binary cubic forms $F$ with integer coefficients, content $1$ and non zero discriminant for which there is a positive number $c$ depending on $F$ such that $N_{F}(m) > c \log(m)^{r/r+2}$ for infinitely many integers $m.$
\end{thm} 

\begin{proof}
    

We utilize a proof of a flavour similar to the one in \ref{BinaryCubic}. In particular, we utilize the syzygy between the covariants of a binary cubic form, adapting it slightly to fit our model. 

Let $P= (s,t)$ be a rational point on $E$ such that $st \neq 0$, then $$F(x,y) = x^3 - 3sx^2y-4Dy^3,$$ with discriminant $D(F) = -432Dt^2$. Let $C$ be the curve $C: F(x,y) = 1/2t$ which is non singular since $st \neq 0.$ We now set $Q  = (-s/t, -1/2t)$ as a rational point on $C$. Taking $Q$ as origin, we realize that $C$ takes the form of an elliptic curve. Now, utilizing the syzygy for the covariants of a binary cubic form, we get $$4 H(x, y)^3=G(x, y)^2+27 D F(x, y)^2.$$ Modifying this equation slightly, we get $$(4G)^2 = (4H)^3 + (432t)^2 DF^2$$ where $$H(x, y) = 9(s^2 x^2 + 4D xy - 4s Dy^2)$$ and $$G(x, y) = 54 \left((s^3 + 2D)x^3 - 6sD x^2 y + 12s^2 D x y^2 + 8D^2 y^3 \right).$$
Now that we have our syzygy sorted, we wish to utilize it somehow to get a possible relation between our curve $C$ and an elliptic curve $E$ using the covariants. This can be achieved by cleverly defining $C: F(x,y) = z^3/2t$ and $E: zy^2 = x^3 + Dz^3$. We now define the mapping  $$ \lambda: C \mapsto E \textrm{ where } \lambda([x, y, z]) = \left(\left[ z H(x, y) / 9, G(x, y) / 54, z^3 \right]\right).$$ The most crucial part of the argument lies here, where we wish to show that $\lambda$ is a non-constant morphism, and is therefore an isogeny with a degree, which allows us to effectively preserve the Mordell-Weil rank. We thus prove the following lemma: 

\begin{lemma}
    The mapping $\lambda$ is a non constant morphism, and therefore is an isogeny relating $C \mapsto E$ while preserving the Mordell-Weil rank. 
\end{lemma}
\begin{proof}
    When $z \neq 0 \text{ or } G(x,y) \neq 0$, we realize that $\lambda$ is regular. Now, when $z = 0 \text{ and } G(x,y) = 0,$  then $F(x,y) =0 $ and by our syzygy, $H(x,y) = 0$. We now realize that $\mathrm{Res}( \frac{H(X, Y)}{9}, \frac{F(X,Y)}{1}) \neq 0$, and hence $\lambda$ is a non constant morphism. This implies that $\lambda$ is an isogeny between the elliptic curves $C$ and $E$ such that the rank over $\mathbb{Q}$ is preserved. 
    \end{proof}
Furthermore, we note that if $Q$ is the origin of the elliptic curve $C$, then $\lambda(Q) = [s, -t, 1]$ is the orgin of the elliptic curve $E$. Now, the last part of proof hinges upon our ability to find binary cubic forms of the same rank as $F(x,y)$, and to prove that there are infinitely many such forms. Let $s = s_1/s_2$ and $t = t_1/t_2$ such that $\gcd(s_1,s_2) = 1$ and $s_2 > 0$ and $\gcd(t_1, t_2) = 1$ with $t_2 > 0$. We now associate $F'$ to be the binary cubic form associated with $F$ via a rank preserving argument, and then show that there can be infinitely many inequivalent forms $F'$. 
\begin{lemma}
    There are infinitely many inequivalent $F'$ such that $\operatorname{rank}F' = \operatorname{rank}F$ with content $1$.
\end{lemma}
\begin{proof}\renewcommand{\qedsymbol}{}

Set $(s,t) = (s_1/s_2, t_1,t_2)$ where $\gcd(s_1,s_2) = \gcd(t_1,t_2) = 1$ and $s_2 \neq 0$ and $t_2 \neq 0.$ We now set $b = s_2/(3,s_2)$ and get $F'(x,y) = b F(x,y).$ Finally, we show that the $\Delta(F') = -432b^4 t^2D, $ then we can simply associate the curve $C_1: F'(x,y) = m_0 z^3$ where $m_0 \neq 0$ and where the content of $F'$ is $1.$ Now utilizing \ref{Silverman's Theorem}, we get $N_{F} (m) > c \log(m)^{r/r+2}$ for infinitely many positive integers $m.$ Lastly, the fact that there are infinitely many inequivalent forms $F'$ can be proved using the fact that there can exist forms $F'$ with discriminants of arbitrarily large absolute value with points associated on $E$, completing our proof.
\end{proof}
\end{proof}

\end{document}
